\documentclass[11pt]{amsart}
\usepackage{amsmath,amssymb,amsthm}
\usepackage[margin=1.1in]{geometry}
\usepackage{hyperref}
\hypersetup{pdftitle={Two of Schnorr's gate elimination bounds, machine-checked in Lean 4},
  pdfauthor={Santiago Gegenschatz}}
\usepackage{fancyvrb}
\usepackage[no-math]{fontspec}
% Lean source is Unicode; Menlo has every symbol the listings use.
\newfontfamily\leanfont{Menlo}[Scale=0.82]

\newtheorem{theorem}{Theorem}
\newtheorem{lemma}[theorem]{Lemma}
\theoremstyle{definition}
\newtheorem{definition}[theorem]{Definition}
\theoremstyle{remark}
\newtheorem{remark}[theorem]{Remark}

\newcommand{\lean}[1]{\texttt{#1}}
\emergencystretch=3em
\DefineVerbatimEnvironment{leancode}{Verbatim}{formatcom=\leanfont,xleftmargin=1.5em}

\title[Two gate elimination bounds in Lean 4]{Two of Schnorr's gate elimination bounds,
machine-checked in Lean 4}
\author{Santiago Gegenschatz}
\address{Independent researcher}
\email{santy@lebenmaster.com}
\date{Draft, 24 November 2026. Not peer reviewed. The Lean code and this note were written by
an AI assistant (Claude Fable 5.2, in Claude Code) under the author's direction.}

\begin{document}

\begin{abstract}
A Boolean circuit whose gates are two-input functions other than XOR and XNOR needs exactly
$3(n-1)$ gates to compute the parity of $n$ inputs. With every two-input gate allowed,
deciding whether at least two of $n$ inputs are on needs at least $2n-3$ gates, and so does
any function that keeps mixing each pair of its free inputs while inputs are fixed one at a
time. Both bounds go back to Schnorr in the 1970s; the general condition in the second is our
variant of his. We give proofs of both checked by the Lean~4 proof assistant. The development is 2{,}395 lines, uses only the Lean core library (no Mathlib), has
no \lean{sorry}, depends only on the three standard axioms, and passes a replay through Lean's
kernel and one independently written proof checker. The theorems are not new. What we offer is the machine check and a way of arranging
the proofs that made them short: gates are never removed and wires are never renumbered; a
gate that has become useless is overwritten in place by a free node. We found no earlier
machine check of these bounds, but we cannot rule one out. The code and the text were written
by an AI assistant under the author's direction. A separate AI referee pass re-ran the checks
and tested the definitions on small circuits. No human has checked the definitions line by
line.
\end{abstract}

\maketitle

\section{What is proved}

A circuit has $n$ input wires. Each gate applies a function $\{0,1\}^2\to\{0,1\}$ to two
earlier wires and adds one new wire. The last wire is the output, and the size is the number
of gates. Of the sixteen two-input functions, call XOR and XNOR \emph{XOR-like}. The other
fourteen, which include the constants and the functions of one argument, form the basis
usually written $U_2$.

\begin{theorem}[Schnorr; Lean: \lean{parity\_exact\_without\_xor}]\label{thm:main}
Every circuit on $n$ inputs with no XOR-like gate that computes
$x_1\oplus\dots\oplus x_n$ has at least $3(n-1)$ gates, and some such circuit has exactly
$3(n-1)$ gates.
\end{theorem}

Here is the statement as Lean accepts it.

\begin{leancode}
theorem parity_exact_without_xor (n : Nat) :
    (∀ c : Circuit, c.inputs = n → (∀ g, g ∈ c.gates → ¬IsXorLike g.op) →
        c.Computes parity → 3 * (n - 1) ≤ c.size) ∧
      ∃ c : Circuit, c.inputs = n ∧ (∀ g, g ∈ c.gates → ¬IsXorLike g.op) ∧
        c.Computes parity ∧ c.size = 3 * (n - 1)
\end{leancode}

The Lean statement has no condition on $n$; for $n=0$ and $n=1$ it says that no gate is needed.

The theorem is attributed to Schnorr. Carmosino, Dang and Jackman \cite[Theorem~18]{CDJ26}
cite \cite{Schnorr74} for the $3(n-1)$ bound over AND, OR and NOT. Theorem~\ref{thm:main} is
the wider form, with all fourteen gates that are not XOR-like. We believe that form is due to
\cite{Schnorr76}, whose title names the equivalence function (the negation of parity), but we
have not confirmed the basis used there. We checked that both references exist and have read
neither. The argument is the standard one (fix one input,
three gates die). We wrote the proof from that outline and did not follow any written proof
line by line.

\begin{theorem}[after Schnorr; Lean: \lean{atLeastTwo\_needs\_two\_gates\_per\_input}]\label{thm:two}
Every circuit on $n$ inputs, with any two-input gates, that decides whether at least two of
its inputs are on has at least $2n-3$ gates.
\end{theorem}

\begin{leancode}
theorem atLeastTwo_needs_two_gates_per_input (c : Circuit)
    (hc : c.Computes atLeastTwo) : 2 * c.inputs - 3 ≤ c.size
\end{leancode}

The names of the second development live in the namespace \lean{PNP.Full}; we print them
without the prefix.

This is an instance of a general statement, Theorem~\ref{thm:hard} below. Its hypothesis is our
own formulation. It is modelled on the condition attributed to Schnorr: every two inputs leave
at least three subfunctions, and this survives the fixing of inputs. Ours asks a little more
(dependence on every free input, and a fixing step all the way down to two free inputs), so it
covers fewer functions. For example, no family makes the three-input multiplexer hard in our
sense, because fixing its selector leaves a function of one input. The abstract of
\cite{Schnorr74} describes a linear lower bound that depends on the set of subfunctions, and we
take that paper to be the source. We have not read it, and we have not checked against any
source that the bound is stated there as $2n-3$ or that this threshold function is among its
examples. For $n=2$ the bound is met by one AND gate (\lean{andPair\_computes\_atLeastTwo}).
We did not formalize an upper bound for larger $n$.

\subsection*{Why check it}
Gate elimination is the method behind the linear lower bounds for general circuits, from
$2n$ to a recent claim of $4n-O(\log n)$ \cite{MS26}, and its proofs are long case analyses. Carmosino, Dang and Jackman
\cite{CDJ26} name as a research direction to ``formalize known gate elimination arguments in
standard proof assistants'' and re-prove the $3(n-1)$ bound on paper with that aim. This note is one
step in that direction: the two smallest gate elimination proofs that have the full shape of
the big ones (pick a gate, fix an input, count the gates that die, repeat).

\subsection*{What we searched}
Eleven web searches in November 2026, in two separate passes, found no proof-assistant check of
these bounds or of any other gate elimination lower bound. Circuit lower bounds proved by other
methods have been formalized: parity for bounded-depth circuits in Lean \cite{ParityAC0}, and
the monotone bound for clique in Isabelle/HOL \cite{Thiemann22}. We searched the Archive of
Formal Proofs and GitHub only through web search. So ``first'' is a claim we do not make.

\section{The model, exactly}

A reader who wants to trust Theorems~\ref{thm:main} and~\ref{thm:two} needs to check only this
section against the files. Everything else is checked by Lean. The listing gathers the
definitions from four files (\lean{Circuit}, \lean{Parity}, \lean{NoXor}, \lean{Threshold}),
in namespaces \lean{PNP}, \lean{PNP.Circuit} and \lean{PNP.Full}, with the comments removed.

\begin{leancode}
structure Gate where
  op : Bool → Bool → Bool
  left : Nat
  right : Nat

structure Circuit where
  inputs : Nat
  gates : List Gate

def size (c : Circuit) : Nat := c.gates.length

def WellFormed (c : Circuit) : Prop :=
  ∀ (i : Nat) (h : i < c.gates.length),
    c.gates[i].left < c.inputs + i ∧ c.gates[i].right < c.inputs + i

def fire (wires : List Bool) (g : Gate) : List Bool :=
  wires ++ [g.op (wires.getD g.left false) (wires.getD g.right false)]

def wires (c : Circuit) (x : List Bool) : List Bool := c.gates.foldl fire x
def eval (c : Circuit) (x : List Bool) : Bool := (c.wires x).getLastD false

def Computes (c : Circuit) (f : List Bool → Bool) : Prop :=
  c.WellFormed ∧ ∀ x : List Bool, x.length = c.inputs → c.eval x = f x

def parity : List Bool → Bool
  | [] => false
  | b :: bs => b != parity bs

def IsXorLike (op : Bool → Bool → Bool) : Prop :=
  (∀ a b, op a b = (a != b)) ∨ (∀ a b, op a b = (a == b))

def ones : List Bool → Nat
  | [] => 0
  | b :: bs => (if b then 1 else 0) + ones bs

def atLeastTwo (x : List Bool) : Bool := decide (2 ≤ ones x)
\end{leancode}

Choices a referee should know about.
\begin{itemize}
\item Wires are numbered $0,\dots,n-1$ for the inputs and $n+t$ for gate $t$. A gate may read
  any earlier wire, so fan-out is unbounded, and it may read the same wire twice.
\item The output is the last wire. For a lower bound this loses nothing: if the output of a
  circuit is an earlier gate, deleting the gates after it gives a smaller circuit in our
  model. With no gates at all the output is the last input, which is the right answer for
  $n=1$.
\item Size counts every gate, including gates that compute a constant, a copy or a negation
  of one argument. Counting them makes the lower bound weaker than the one we prove inside
  (Theorem~\ref{thm:and}), where they are free.
\item Subtraction on \lean{Nat} is truncated, so for $n=0$ the bound reads $0\le$ size.
\end{itemize}

\section{The proof of the $3(n-1)$ bound}

\subsection{Nodes instead of gates}
A gate that is not XOR-like is one of three things. Lean checks this by a table of the sixteen
truth tables (\lean{classify}, \lean{exists\_kind}).

\begin{leancode}
inductive Kind where
  | and (a b c : Bool)   -- ((x != a) && (y != b)) != c
  | copyL (e : Bool)     -- x != e
  | copyR (e : Bool)     -- y != e
  | const (d : Bool)     -- d
\end{leancode}

A \emph{node} is a kind with two wire numbers. The eight \lean{and} kinds are AND with any
choice of negations on the two inputs and the output; this covers AND, OR, NAND, NOR, the
two implications and their two negations. The \emph{cost} of a list of nodes is the number of \lean{and} nodes. Copies,
negations and constants are free.

\begin{theorem}[Lean: \lean{and\_gates\_for\_xor}]\label{thm:and}
Let $N$ be a well-formed list of nodes on $n$ inputs. Let $P$ be a nonempty set of inputs in which the
positions in a list $F$ vary freely and every other position is held at a fixed value. If the
last wire of $N$ equals $\bigoplus_{i\in F}x_i$ or its negation on all of $P$, then
$N$ has at least $3(|F|-1)$ \lean{and} nodes.
\end{theorem}

Theorem~\ref{thm:main} follows with $F=\{0,\dots,n-1\}$ and $P$ all inputs of length $n$,
because cost is at most size.

\subsection{Never delete, overwrite}
On paper, gate elimination removes gates and rewires the circuit. In a proof assistant that
means renumbering every wire after each removal, and most of the work goes into the
renumbering. We avoid it. The list of nodes keeps its length and every wire keeps its number
for the whole proof. Two things change instead.
\begin{itemize}
\item \emph{Fixing an input} does not touch the circuit. It shrinks the set $P$ of inputs on
  which the circuit is required to be correct (\lean{fixInput}, \lean{Cube.restrict}).
\item \emph{Killing a gate} overwrites one \lean{and} node by a free node. Call an \lean{and}
  node \emph{dead on $P$} if one of the two wires it reads is constant on $P$. A dead node
  agrees on $P$ with a copy of its other argument or with a constant, so it can be overwritten
  with every wire value on $P$ unchanged and the cost one lower (\lean{kill\_one}). Three
  distinct dead nodes give cost three lower (\lean{kill\_three}).
\end{itemize}

Because copies are free, the proof needs to see through them. \lean{Alias n N s w e} says
that wire $w$ is wire $s$ passed through a chain of copy nodes, negated in total $e$ times
modulo two. Three lemmas about it carry the argument.
\begin{itemize}
\item \lean{exists\_root}: every wire is an alias of an input, of an \lean{and} node, or of
  a constant node.
\item \lean{flow}: take two runs whose inputs differ only at wire $s$ and whose node lists
  differ only at the node of wire $s$. If every other \lean{and} node that reads an alias of
  $s$ takes the same value in both runs, then so does every wire that is not an alias of $s$.
  (In the induction below $s$ is an input and the two lists are equal;
  \lean{reader\_or\_useless} uses the general form.)
\item \lean{reader\_or\_useless}: for an \lean{and} node $g$, either another \lean{and} node
  reads an alias of $g$, or the output is an alias of $g$, or some list of the same length with one \lean{and}
  node fewer has the same output on every input.
\end{itemize}

\subsection{The induction}
The proof of Theorem~\ref{thm:and} is by strong induction on the cost, for all $F$ and $P$ at
once. If $|F|\le 1$ there is nothing to prove. Whenever some \lean{and} node can be overwritten
with the output unchanged on $P$, the induction hypothesis for the same $F$ and $P$ finishes.
We call that an \emph{easy exit}. Otherwise:

\begin{enumerate}
\item There is an \lean{and} node; without one the output would be an alias of one input or a
  constant, and the XOR of two free inputs is neither. Let $g$ be the first one. Both wires it
  reads lie before every \lean{and} node, so each is a constant (easy exit) or an alias of an
  input. Say they are aliases of inputs $i$ and $j$.
\item If $i=j$, then $g$ computes a function of one input and is overwritten by a copy or a
  constant (easy exit). If $i$ or $j$ is not in $F$, it is constant on $P$ and $g$ is dead
  (easy exit). So $i\ne j$ and both are free.
\item A second \lean{and} node $h\ne g$ reads an alias of $i$. If none did, set $x_j$ to the
  value that forces $g$ and flip $x_i$: the output of $g$ stays put, so by \lean{flow} nothing
  but aliases of $i$ would change. The output is not an alias of $i$ (it depends on $j$ as
  well), so the output would not change. But the output is the XOR of the free inputs and $i$
  is free.
\item Fix $x_i$ to the value that forces the output of $g$. Now $P$ is smaller, $F$ has lost
  $i$, and the last wire still computes the XOR of the remaining free inputs up to negation.
  Nodes $g$ and $h$ each read a stuck wire and are dead.
\item The output of $g$ is now constant. By \lean{reader\_or\_useless}, either an \lean{and}
  node $p$ reads an alias of $g$, or the output is an alias of $g$ and so constant
  (impossible, since $j$ is still free), or there is an easy exit. Node $p$ comes after $g$ and is dead.
\item If $p\ne h$, the nodes $g,h,p$ are three distinct dead nodes. If $p=h$, then $h$ reads
  an alias of $i$ and an alias of $g$, both constant now, so the output of $h$ is constant
  too. Apply \lean{reader\_or\_useless} to $h$: some \lean{and} node $q$ after $h$ reads an
  alias of $h$ and is dead, and $g,h,q$ are distinct.
\item Overwrite the three dead nodes. The cost drops by three, $|F|$ drops by one, and the
  induction hypothesis gives $3(|F|-2)\le \text{cost}-3$.
\end{enumerate}

In Lean this is one proof of about 190 lines (\lean{Schnorr.lean}).

\subsection{The matching circuit}
Write $x\oplus y=(x\wedge\neg y)\vee(\neg x\wedge y)$: three gates, none XOR-like. Chaining
$n-1$ of these gives a circuit of size $3(n-1)$, and Lean checks that it is well formed, has
no XOR-like gate and computes parity (\lean{andChain}, \lean{andChain\_computes\_parity}).

\section{The $2n-3$ bound}

Now every two-input function is a gate. The same overwrite-in-place setup is built once more
with a node that is an arbitrary gate, a copy or a constant (\lean{FullBasis.lean},
\lean{FullAlias.lean}). The one fact about gates that the proof uses is that a gate with one
wire stuck is a function of its other wire, so it is a copy or a constant. That holds for XOR
as well as for AND, which is why fixing an input here kills two gates and not three.

Let $P$ be a set of allowed inputs with free positions $F$, as before. Say that $f$
\emph{mixes} positions $i$ and $j$ on $P$ if no single bit computed from $x_i$ and $x_j$ is
enough for $f$:

\begin{leancode}
def Mixes (P : List Bool → Prop) (f : List Bool → Bool) (i j : Nat) : Prop :=
  ∀ b : Bool → Bool → Bool, ∃ x x', P x ∧ P x' ∧
    (∀ w, w ≠ i → w ≠ j → x.getD w false = x'.getD w false) ∧
    b (x.getD i false) (x.getD j false) = b (x'.getD i false) (x'.getD j false) ∧
    f x ≠ f x'
\end{leancode}

This says the same as the usual condition: the four ways of fixing $x_i$ and $x_j$ leave at
least three different functions of the other inputs. (If they left at most two, the bit that
says which of the two would be enough; and a bit that is enough leaves at most two.) We did not
formalize this equivalence.

A family $Q$ of pairs $(F,P)$ is \emph{hard for $f$} (\lean{Hard}) if on each of its members
$f$ depends on every free position when $|F|\ge 2$, mixes every two free positions when
$|F|\ge 3$, and, when $|F|\ge 3$, each free position can be fixed to some value so that the
restricted pair is again in $Q$.

\begin{theorem}[Lean: \lean{hard\_needs\_two\_per\_input}]\label{thm:hard}
If $Q$ is hard for $f$ and contains the pair of all positions and all inputs of length $n$,
then every circuit on $n$ inputs that computes $f$ has at least $2n-3$ gates.
\end{theorem}

The proof is by strong induction on the number of gates. Easy exits are as in Section~3. If
$|F|=2$ the claim is one gate, and a circuit with no gate outputs a copy of one input or a
constant. Otherwise let $g$ be the first gate; it reads aliases of two free inputs $i\ne j$.
\begin{enumerate}
\item The output is not an alias of $i$ or of $j$: the bit $x_i$ alone, or $x_j$ alone, would
  then be enough for $f$, and $f$ mixes $i$ and $j$.
\item Some gate $h\ne g$ reads an alias of $i$ or of $j$. If none did, apply \lean{Mixes} to
  the bit that $g$ puts out: two inputs agree outside $i,j$ and on that bit. By \lean{flow},
  now with the two sources $i$ and $j$, every wire that is not an alias of $i$ or $j$ agrees on
  the two inputs, so the output does. But $f$ differs on them.
\item Say $h$ reads an alias of $s\in\{i,j\}$. Fix $x_s$ to the value hardness provides. Both
  $g$ and $h$ read a stuck wire and are dead. Overwrite them: the cost drops by two and
  $|F|$ by one.
\end{enumerate}

For Theorem~\ref{thm:two}, $Q$ is the family of pairs in which the all-off input is allowed,
so every fixed input is off. With $x_i=x_j=0$ the other inputs must supply two ones, with
exactly one of them on they must supply one, and with both on nothing is needed. A third free
input tells these three apart (\lean{Threshold.lean}).

\section{How it was checked}

\begin{itemize}
\item Lean 4.36.0, core library only. The whole project builds in a few seconds on a laptop.
\item The Lean source and the source of this note are public:
  \url{https://github.com/santy-gegenschatz/gate-elimination-lean}. The file
  \lean{Axioms.lean} there prints the axioms of the three theorems.
\item \lean{\#print axioms} for Theorems~\ref{thm:main}, \ref{thm:two} and~\ref{thm:hard} reports \lean{propext},
  \lean{Classical.choice} and \lean{Quot.sound}, and nothing else.
\item A text scan of the source finds no \lean{sorry}, \lean{axiom}, \lean{native\_decide},
  \lean{unsafe}, \lean{implemented\_by}, \lean{extern}, \lean{opaque} or \lean{partial}.
\item \lean{leanchecker}, which replays every compiled module through Lean's kernel, accepts
  the whole library.
\item \lean{nanoda}, a proof checker written independently in Rust, checked the 2{,}694
  declarations the three theorems depend on, exported with \lean{lean4export}, with no errors.
  These are the checks that were applied by others to the Lean proof of the generalized
  Riemann hypothesis announced in November 2026 \cite{Rubinstein26}; we copied the list from there.
\item For the first bound the matching circuit is a control on the statement: a lower bound stated too strongly
  would contradict it inside Lean.
\item The files, in dependency order: \lean{Circuit}, \lean{Reads}, \lean{Parity},
  \lean{AndBasis}, \lean{Alias}, \lean{Kill}, \lean{Cube}, \lean{Schnorr}, \lean{NoXor},
  \lean{Tight} for the first bound, then \lean{FullBasis}, \lean{FullAlias}, \lean{TwoN},
  \lean{Threshold} for the second; 2{,}395 lines in total. \lean{Reads} and \lean{Parity} hold earlier folklore
  results (with XOR gates allowed, parity needs exactly $n-1$ gates) and are used here only
  for the definition of \lean{parity} and a few list lemmas.
\end{itemize}

\section{Limits of this note}

\begin{itemize}
\item The theorems are about fifty years old. Nothing here moves any lower bound.
\item A machine check proves that the Lean statement follows from the Lean definitions. That
  the definitions say what Section~2 claims has been checked by the AI assistant that wrote
  them and by one separate AI referee pass, which was asked to break the note, re-ran every
  check in Section~5 and ran the definitions on small circuits. No
  human has checked them line by line.
\item Theorem~\ref{thm:hard} is stated for our own class of hard functions. We have not compared
  it clause by clause with a published statement of Schnorr's class.
\item The search for earlier formalizations was a web search, not a survey.
\item The two developments share their shape but not their code: the nodes of the first have
  no XOR gate, the nodes of the second have arbitrary gates. Unifying them is left for the next
  bound.
\item Both bounds are far from the frontier. For the full basis a bound of $4n-O(\log n)$ was
  claimed in November 2026 \cite{MS26}; it is not yet refereed. The plan is to carry the same
  setup up the line of gate elimination bounds.
\end{itemize}

\begin{thebibliography}{9}
\bibitem{CDJ26} M.~Carmosino, N.~Dang, T.~Jackman, \emph{Convergent gate elimination and
  constructive circuit lower bounds}, arXiv:2602.17942, 2026.
\bibitem{MS26} O.~S.~Mishra, R.~Santhanam, \emph{An explicit $4n-O(\log n)$ size lower bound
  against general Boolean circuits}, ECCC report TR26-260, revision~1, 3 November 2026.
\bibitem{ParityAC0} S.~Wesonga, \emph{Formalizing PARITY circuit lower bounds in Lean},
  arXiv:2609.24188, 2026.
\bibitem{Rubinstein26} M.~Rubinstein, \emph{Alp\"oge/Claude proof of the Generalized Riemann
  Hypothesis}, web page, November 2026,
  \url{https://www.math.uwaterloo.ca/~mrubinst/alpoge_claude_grh/}.
\bibitem{Schnorr74} C.~P.~Schnorr, \emph{Zwei lineare untere Schranken f\"ur die Komplexit\"at
  Boolescher Funktionen}, Computing \textbf{13} (1974), no.~2, 155--171.
\bibitem{Schnorr76} C.~P.~Schnorr, \emph{The combinational complexity of equivalence},
  Theoretical Computer Science \textbf{1} (1976), no.~4, 289--295.
\bibitem{Thiemann22} R.~Thiemann, \emph{Clique is not solvable by monotone circuits of
  polynomial size}, Archive of Formal Proofs, May 2022,
  \url{https://isa-afp.org/entries/Clique_and_Monotone_Circuits.html}.
\end{thebibliography}

\end{document}
